\section{Resultados con los datos del HUCA}
\label{chap:comparacion_real}

Todos los resultados de este apartado utilizan las etiquetas de los pacientes descritos en la sección~\ref{sec:procedencia_huca}.

\subsection{Transferencia desde el simulador}

La CNN entrenada con el simulador alcanza su mejor BA sobre el HUCA con \(K=4\). Al aumentar los ejemplos mejora en los datos generados, pero cae hasta 0,498 sobre los ECG reales con \(K=32\). Aprender mejor las reglas del simulador no se traduce, por tanto, en reconocer mejor a los pacientes del HUCA. La Tabla~\ref{tab:domain_adaptation_summary} compara ese mejor punto con CORAL y MMD.

\begin{table}[H]
  \centering
  \small
  \begin{tabularx}{0.95\textwidth}{Xrrrrr}
    \toprule
    Configuración & \(K\) & BA & F1 & Sens. & Esp. \\
    \midrule
    CNN base, mejor global & 4 & \(0,551 \pm 0,034\) & 0,506 & 0,670 & 0,433 \\
    CORAL & 16 & \(0,545 \pm 0,011\) & 0,501 & 0,670 & 0,421 \\
    MMD & 32 & \(0,546 \pm 0,011\) & 0,502 & 0,678 & 0,413 \\
    \bottomrule
  \end{tabularx}
  \caption{Puntos principales de la CNN sobre los datos del HUCA. CORAL y MMD se comparan con sus bases del mismo valor de \(K\).}
  \label{tab:domain_adaptation_summary}
  \resultstablenote
\end{table}

MMD eleva la BA de 0,498 a 0,546 con \(K=32\) y CORAL, de 0,520 a 0,545 con \(K=16\). La mejora es limitada y el rendimiento sigue cerca del azar.

Con demostraciones del simulador y selección balanceada, Gemma~4 obtiene su mejor BA, \(0,514 \pm 0,008\), en \(K=16\), con sensibilidad de 0,063 y especificidad de 0,965. La mayoría de los positivos no se detecta, lo que muestra que estos ejemplos tampoco permiten trasladar a los ECG del HUCA el reconocimiento conseguido en el simulador. En \(K=32\), la BA baja a 0,512, la sensibilidad sube a 0,072 y la especificidad baja a 0,952.

En conjunto, ni la CNN ni Gemma~4 consiguen generalizar adecuadamente a los registros reales con los ejemplos del simulador. La Figura~\ref{fig:transfer_summary} resume las cuatro variantes de adaptación y Gemma~4 con demostraciones del simulador.

\begin{figure}[H]
  \centering
  \safeincludegraphics[width=0.95\textwidth]{results/transfer_controls_synthetic_to_real_summary.pdf}
  \caption{Transferencia desde el simulador a los ECG clínicos. A la izquierda se comparan las variantes de adaptación de la CNN. A la derecha se muestra Gemma~4 con demostraciones del simulador.}
  \label{fig:transfer_summary}
\end{figure}

\subsection{Aprendizaje con ejemplos clínicos}
\label{sec:resultados_clinicos}

Conforme recibe más ejemplos del conjunto cedido por el HUCA, Gemma~4 mejora la clasificación de pacientes que no forman parte de las demostraciones de su consulta. Esta generalización contrasta con la CNN entrenada con esos datos, que permanece cerca del azar y no mejora de forma continua al aumentar los ejemplos. La Tabla~\ref{tab:comparacion_clinica_cnn} recoge las mejores configuraciones de ambos modelos.

\IfFileExists{generated/clinical_cnn_comparison_table.tex}{\begin{table}[H]
  \centering
  \small
  \begin{tabularx}{\textwidth}{Xrrrrr}
    \toprule
    Método & \(K\) & BA & F1 & Sens. & Esp. \\
    \midrule
    CNN estándar & 32 & \(0,532 \pm 0,002\) & \(0,469 \pm 0,007\) & \(0,580 \pm 0,018\) & \(0,483 \pm 0,013\) \\
    CNN balanceada & 16 & \(0,546 \pm 0,017\) & \(0,478 \pm 0,017\) & \(0,580 \pm 0,028\) & \(0,511 \pm 0,033\) \\
    Gemma 4 estándar & 32 & \(0,743 \pm 0,011\) & \(0,675 \pm 0,014\) & \(0,690 \pm 0,014\) & \(0,796 \pm 0,008\) \\
    Gemma 4 balanceado & 32 & \(0,761 \pm 0,011\) & \(0,699 \pm 0,013\) & \(0,741 \pm 0,014\) & \(0,781 \pm 0,008\) \\
    \bottomrule
  \end{tabularx}
  \caption{Mejores configuraciones de la CNN entrenada con el HUCA y de Gemma~4 con demostraciones de ese mismo conjunto.}
  \label{tab:comparacion_clinica_cnn}
  \resultstablenote
\end{table}
}{}

Sin demostraciones, Gemma~4 obtiene una BA de 0,499 y no identifica ningún positivo del HUCA. Al igual que en el simulador, su configuración \textit{zero-shot} no consigue resolver la tarea con las instrucciones utilizadas. La mejor CNN supera ligeramente a Gemma~4 con \(K=2\), 0,519 frente a 0,492, y con \(K=4\), 0,520 frente a 0,503. Desde \(K=8\) cambia el orden y Gemma~4 con selección balanceada alcanza 0,661 frente a 0,507 de la mejor CNN. Con \(K=16\), las mejores BA son 0,710 y 0,546, y con \(K=32\), 0,761 y 0,535.

\begin{figure}[H]
  \centering
  \safeincludegraphics[width=0.98\textwidth]{results/clinical_cnn_vs_gemma_performance.pdf}
  \caption{BA y F1 de la CNN y Gemma~4 al aprender con ejemplos del HUCA. Las barras representan la desviación estándar de tres semillas.}
  \label{fig:clinical_cnn_gemma_performance}
\end{figure}

La Figura~\ref{fig:clinical_cnn_gemma_performance} muestra la separación en BA y F1. La diferencia entre los mejores puntos es 0,216 de BA, y con \(K=32\) fijo es 0,226.

Las matrices de la Figura~\ref{fig:clinical_cnn_gemma_confusions} muestran que Gemma~4 mejora la clasificación de ambas clases. Con ejemplos del HUCA reconoce positivos y negativos de pacientes que no forman parte de las demostraciones de su consulta. Esta generalización dentro del conjunto contrasta con el resultado obtenido sin ejemplos o con ejemplos del simulador, pero no demuestra todavía el mismo rendimiento en otros centros. Cada matriz reúne tres evaluaciones de los mismos pacientes.

\begin{figure}[H]
  \centering
  \safeincludegraphics[width=0.90\textwidth]{results/clinical_cnn_vs_gemma_confusions.pdf}
  \caption{Matrices de confusión normalizadas por clase real para la CNN entrenada con selección balanceada de pacientes del HUCA en \(K=16\) y Gemma~4 con demostraciones del HUCA balanceadas en \(K=32\).}
  \label{fig:clinical_cnn_gemma_confusions}
\end{figure}

En esta comparación, Gemma~4 detecta más positivos y genera menos falsas alertas que la CNN, lo que apoya investigar su uso para priorizar revisiones en el triaje. Los positivos que omite impiden plantearlo como filtro de descarte y hacen necesario mantener la evaluación clínica completa.

\subsection{Comparación entre el simulador y el HUCA}

Las Figuras~\ref{fig:domain_gap_cnn} y~\ref{fig:domain_gap_icl} permiten comparar las curvas completas. La CNN pasa de un máximo de 0,885 en el simulador a 0,551 sobre el HUCA. Gemma~4 pasa de 0,904 en el simulador a 0,761 con consultas y ejemplos clínicos. Con ejemplos del simulador y consulta clínica, su máximo es 0,514.

\begin{figure}[H]
  \centering
  \safeincludegraphics[width=0.92\textwidth]{results/domain_gap_cnn_ba_by_k.pdf}
  \caption{Exactitud balanceada de la CNN entrenada en el simulador y evaluada con imágenes del simulador o clínicas, según el número de ejemplos.}
  \label{fig:domain_gap_cnn}
\end{figure}

\begin{figure}[H]
  \centering
  \safeincludegraphics[width=0.94\textwidth]{results/domain_gap_icl_context_ba_by_k.pdf}
  \caption{Exactitud balanceada de Gemma~4 según el número de ejemplos y el origen de las demostraciones y de la consulta.}
  \label{fig:domain_gap_icl}
\end{figure}
